\chapter{Diophantine Equations}\label{diophantine-equations}
\chapterstatus{draft}

A \emph{Diophantine equation} is a polynomial equation with integer coefficients of which one seeks the integer, or the rational, solutions. The subject has no general method; Hilbert's tenth problem,
discussed at the end of the chapter, shows that there cannot be one. What it has instead is a stock of ideas, each of which settles whole families of equations: congruences and $p$-adic numbers show that an
equation has no solutions; Bézout's identity and geometry (lines through a known point) produce all of them; Fermat's \emph{infinite descent} turns a solution into a smaller one; and factorisation in rings
larger than $\ZZ$, such as the Gaussian and Eisenstein integers, turns a sum into a product. This chapter presents these ideas on the classical examples: linear equations and the coin problem, Pythagorean
triples and conics, the equations $x^4 \pm y^4 = z^2$ and congruent numbers, sums of three squares, Markov's equation, Mordell's equation, and Fermat's last theorem for the exponents $3$ and $4$ together with
Sophie Germain's theorem and the polynomial analogue. Pell's equation was treated in Chapter~\ref{continued-fractions}, the Hasse--Minkowski theorem in Chapter~\ref{p-adic-numbers}, and Waring's problem in
Chapter~\ref{additive-number-theory}.

\noindent\textbf{Prerequisites.} Chapter~\ref{number-theory} (Elementary Number Theory), Chapter~\ref{continued-fractions} (Continued Fractions and Diophantine Approximation), Chapter~\ref{additive-number-theory} (Additive Number Theory) and Chapter~\ref{p-adic-numbers} (p-adic Numbers).

\section{Linear equations and the coin problem}\label{diophantine-equations:section-linear}

\begin{notation}\label{diophantine-equations:notation-diophantine}
In this chapter $\gcd(a, b)$ is the greatest common divisor of integers not both zero, taken positive, and $a \mid b$ means that $a$ divides $b$. An \emph{integer solution} of an equation
$f(x_1, \dots, x_n) = 0$ with $f \in \ZZ[x_1, \dots, x_n]$ is a point of $\ZZ^n$ where $f$ vanishes, a \emph{rational solution} a point of $\QQ^n$. A solution is \emph{trivial} when it is the obvious one an
equation of the given shape always has (for instance $x = y = z = 0$ for a homogeneous equation), and \emph{primitive} when its coordinates have greatest common divisor $1$.
\end{notation}

The simplest Diophantine equations are linear, and Bézout's theorem (Theorem~\ref{number-systems:theorem-bezout}) solves them completely.

\begin{proposition}\label{diophantine-equations:proposition-linear}
Let $a, b, c \in \ZZ$ with $a, b$ not both zero, and $d = \gcd(a, b)$. The equation $ax + by = c$ has an integer solution if and only if $d \mid c$. If $(x_0, y_0)$ is one solution, all solutions are
\[
x = x_0 + \frac bd k, \qquad y = y_0 - \frac ad k, \qquad k \in \ZZ.
\]
\end{proposition}

\begin{proof}
\emph{Necessity.} $d$ divides $ax + by$ for all integers $x$, $y$.

\emph{Sufficiency.} By Bézout there are $u$, $v$ with $au + bv = d$; if $c = de$, then $(ue, ve)$ is a solution.

\emph{All solutions.} The displayed pairs are solutions, since $a\frac bd k - b\frac ad k = 0$. Conversely, if $(x, y)$ is a solution, then $a(x - x_0) = -b(y - y_0)$, so
$\frac ad(x - x_0) = -\frac bd(y - y_0)$. Say $b \neq 0$. As $\gcd(\frac ad, \frac bd) = 1$, Euclid's lemma (Lemma~\ref{number-systems:lemma-euclid}) gives $\frac bd \mid x - x_0$, say $x - x_0 = \frac bd k$,
and then $y - y_0 = -\frac ad k$. If $b = 0$, then $a \neq 0$, $x = x_0$ is forced and $d = |a|$, so $y$ is arbitrary, which is the formula again.
\end{proof}

\begin{example}\label{diophantine-equations:example-linear}
$\gcd(12, 42) = 6$, so $12x + 42y = 15$ has no solution (it would make $15$ divisible by $6$), while $12x + 42y = 30$ has the solution $(-1, 1)$ and hence the solutions $(-1 + 7k, 1 - 2k)$. With more unknowns
the same holds: $a_1x_1 + \dots + a_nx_n = c$ is solvable if and only if $\gcd(a_1, \dots, a_n) \mid c$, by induction on $n$, since $a_1x_1 + a_2x_2$ runs through the multiples of $\gcd(a_1, a_2)$.
\end{example}

A different question arises when the unknowns must be \emph{nonnegative}: which amounts can be paid exactly with coins of values $a$ and $b$? Not all, since the amounts $1, \dots, \min(a, b) - 1$ cannot; but all
large amounts can, and the answer to "how large" is a theorem of Sylvester.

\begin{theorem}[Sylvester]\label{diophantine-equations:theorem-frobenius}
Let $a, b \geq 1$ be coprime integers, and call $n \in \ZZ$ \emph{representable} if $n = ax + by$ with integers $x, y \geq 0$. Then:
\begin{enumerate}
\item every $n > ab - a - b$ is representable, and $ab - a - b$ is not;
\item of the integers $0, 1, \dots, ab - a - b$, exactly half, namely $\frac{(a - 1)(b - 1)}{2}$, are not representable.
\end{enumerate}
\end{theorem}

\begin{proof}
\emph{Step 1: a normal form.} Every integer $n$ can be written uniquely as $n = ax + by$ with $x, y \in \ZZ$ and $0 \leq x \leq b - 1$: by Proposition~\ref{diophantine-equations:proposition-linear} the
solutions of $ax + by = n$ are $(x_0 + bk, y_0 - ak)$, and exactly one $k$ puts $x_0 + bk$ in $[0, b - 1]$. If $n = ax^{\prime} + by^{\prime}$ with $x^{\prime}, y^{\prime} \geq 0$, then replacing $x^{\prime}$ by its
remainder modulo $b$ increases $y^{\prime}$; so \emph{$n$ is representable if and only if $y \geq 0$ in its normal form}.

\emph{Step 2: the bound.} The normal form of $N = ab - a - b$ is $a(b - 1) + b(-1)$, with $y = -1$, so $N$ is not representable. If $n > N$ has normal form $(x, y)$, then
$by = n - ax \geq n - a(b - 1) > N - a(b - 1) = -b$, so $y > -1$ and $n$ is representable.

\emph{Step 3: the count.} If $n$ has normal form $(x, y)$, then $N - n = a(b - 1 - x) + b(-1 - y)$ is the normal form of $N - n$, as $0 \leq b - 1 - x \leq b - 1$. Exactly one of $y$ and $-1 - y$ is $\geq 0$, so
exactly one of $n$, $N - n$ is representable. So $n \mapsto N - n$ is an involution of $\{0, \dots, N\}$ exchanging representable and non-representable integers (in particular it has no fixed point), and the
$N + 1 = (a - 1)(b - 1)$ integers split into two equal halves.
\end{proof}

\begin{example}\label{diophantine-equations:example-frobenius}
With stamps of $3$ and $5$ cents one can pay every amount from $8$ cents on, but not $7 = 15 - 3 - 5$; the non-payable amounts are $1, 2, 4, 7$, four of them, as $\frac{2 \cdot 4}{2} = 4$. The largest
non-representable integer for a set of coprime generators is called its \emph{Frobenius number}. For three or more generators no formula like $ab - a - b$ is known; for $\{6, 9, 20\}$ it is $43$
(Exercise~\ref{diophantine-equations:exercise-frobenius-three}).
\end{example}

\begin{counterexample}\label{diophantine-equations:counterexample-frobenius}
The coprimality hypothesis is necessary: if $\gcd(a, b) = d > 1$, no integer prime to $d$ is of the form $ax + by$, so there are arbitrarily large non-representable integers.
\end{counterexample}

\section{Pythagorean triples and conics}\label{diophantine-equations:section-conics}

\begin{definition}\label{diophantine-equations:definition-pythagorean}
A \emph{Pythagorean triple} is a triple $(x, y, z)$ of positive integers with $x^2 + y^2 = z^2$; it is \emph{primitive} if $\gcd(x, y, z) = 1$.
\end{definition}

Every triple is a multiple of a primitive one, and in a primitive triple $x$, $y$, $z$ are pairwise coprime (a prime dividing two of them divides the square of the third). Exactly one of $x$, $y$ is even: not
both, by coprimality, and not neither, since then $z^2 \equiv 2 \pmod 4$, while squares are $0$ or $1$ modulo $4$. The basic tool is the following lemma, used repeatedly in this chapter.

\begin{lemma}\label{diophantine-equations:lemma-coprime-powers}
Let $R$ be a unique factorisation domain, $k \geq 2$, and $u_1, \dots, u_r \in R$ nonzero and pairwise coprime (no prime divides two of them). If $u_1 \cdots u_r = w^k$, then each $u_i = \varepsilon_iw_i^k$ with
$\varepsilon_i$ a unit and $w_i \in R$. The same holds if $u_1 \cdots u_r = \varepsilon w^k$ with $\varepsilon$ a unit. For $R = \ZZ$ and positive $u_i$, each $u_i$ is a $k$-th power.
\end{lemma}

\begin{proof}
For a prime $\pi$ of $R$, the exponent of $\pi$ in $\varepsilon w^k$ is a multiple of $k$; it is the exponent of $\pi$ in the unique $u_i$ that $\pi$ divides. So in the factorisation of $u_i$ every exponent is a multiple
of $k$, and $u_i = \varepsilon_i\prod\pi^{k e_\pi}$. In $\ZZ$ the units are $\pm1$, and a positive $u_i$ has $\varepsilon_i = 1$ when the primes are taken positive.
\end{proof}

\begin{theorem}\label{diophantine-equations:theorem-pythagorean}
The primitive Pythagorean triples with $y$ even are exactly
\[
x = m^2 - n^2, \qquad y = 2mn, \qquad z = m^2 + n^2,
\]
with $m > n > 0$ coprime integers of opposite parity; distinct pairs $(m, n)$ give distinct triples.
\end{theorem}

\begin{proof}
\emph{These are triples.} $(m^2 - n^2)^2 + (2mn)^2 = (m^2 + n^2)^2$. A prime dividing $x$ and $z$ divides $2m^2$ and $2n^2$; it is odd since $z$ is odd, so it divides $m$ and $n$, which is excluded.

\emph{All triples arise.} Let $(x, y, z)$ be primitive with $y$ even, so $x$, $z$ odd. Then $y^2 = (z - x)(z + x)$, and $\frac{z - x}{2}$, $\frac{z + x}{2}$ are positive integers whose sum $z$ and
difference $x$ are coprime, so they are coprime. Their product is $(y/2)^2$, so by Lemma~\ref{diophantine-equations:lemma-coprime-powers} $\frac{z + x}{2} = m^2$ and $\frac{z - x}{2} = n^2$ with $m, n > 0$.
Then $z = m^2 + n^2$, $x = m^2 - n^2$ (so $m > n$), $y = 2mn$; $m$, $n$ are coprime because $m^2$, $n^2$ are, and of opposite parity because $z$ is odd.

\emph{Injectivity.} $m^2 = \frac{z + x}{2}$ and $n^2 = \frac{z - x}{2}$ are determined by the triple.
\end{proof}

\begin{example}\label{diophantine-equations:example-pythagorean}
$(m, n) = (2, 1), (3, 2), (4, 1), (4, 3), (5, 2)$ give $(3, 4, 5)$, $(5, 12, 13)$, $(15, 8, 17)$, $(7, 24, 25)$, $(21, 20, 29)$. The pair $(3, 1)$ is excluded (same parity): it gives $(8, 6, 10)$, which is not
primitive. Since $p = m^2 + n^2$ for a prime $p \equiv 1 \pmod 4$ (Theorem~\ref{number-theory:theorem-fermat-two-squares}), every such prime is the hypotenuse of a primitive right triangle, which agrees with
the factorisation $z = (m + ni)(m - ni)$ in $\ZZ[i]$ (Section~\ref{number-theory:section-gaussian}).
\end{example}

The same result has a geometric proof which applies to every conic. A primitive triple gives the rational point $(x/z, y/z)$ on the unit circle $X^2 + Y^2 = 1$, and conversely. The circle has the rational
point $(-1, 0)$; a line through it with rational slope $t$ meets the circle in a second point, which is rational because it is the other root of a quadratic equation with rational coefficients of which one root
is rational.

\begin{proposition}\label{diophantine-equations:proposition-circle}
The rational points of $X^2 + Y^2 = 1$ are $(-1, 0)$ and the points
\[
\Bigl(\frac{1 - t^2}{1 + t^2}, \frac{2t}{1 + t^2}\Bigr), \qquad t \in \QQ,
\]
and distinct $t$ give distinct points.
\end{proposition}

\begin{proof}
The line $Y = t(X + 1)$ meets the circle where $X^2 + t^2(X + 1)^2 = 1$, that is, $(X + 1)\bigl((1 + t^2)X - (1 - t^2)\bigr) = 0$. Besides $X = -1$ the root is $X = \frac{1 - t^2}{1 + t^2}$, and then
$Y = t(X + 1) = \frac{2t}{1 + t^2}$. Conversely, a rational point $(X, Y) \neq (-1, 0)$ has $X \neq -1$ (the only point with $X = -1$ is $(-1, 0)$) and lies on the line through $(-1, 0)$ of slope
$t = Y/(X + 1) \in \QQ$, so it is the point with parameter $t$. Distinct slopes give distinct lines through $(-1, 0)$, which meet the circle in distinct second points.
\end{proof}

Writing $t = n/m$ in lowest terms gives $(X, Y) = \bigl(\frac{m^2 - n^2}{m^2 + n^2}, \frac{2mn}{m^2 + n^2}\bigr)$, which is Theorem~\ref{diophantine-equations:theorem-pythagorean} up to the bookkeeping of
common factors. The argument works for any conic with one known rational point.

\begin{proposition}\label{diophantine-equations:proposition-conic}
Let $f \in \QQ[X, Y]$ have degree $2$ and be irreducible over $\CC$, let $C = \{f = 0\}$, and let $P \in C(\QQ) = C \cap \QQ^2$. Then every line through $P$ defined over $\QQ$ meets $C$ in at most one further
point, which is rational, and every point of $C(\QQ) \setminus \{P\}$ arises from exactly one such line. In particular, if $C(\QQ)$ is nonempty, it is infinite.
\end{proposition}

\begin{proof}
Write $P = (p, q)$ and parametrise a line through $P$ with rational direction $(u, v) \neq (0, 0)$ as $(p + su, q + sv)$. Then $g(s) = f(p + su, q + sv)$ is a polynomial in $s$ of degree at most $2$ with
rational coefficients and $g(0) = 0$, say $g(s) = s(\alpha + \beta s)$. It is not identically zero: otherwise the line would lie in $C$, and the linear polynomial defining it would divide $f$ over $\CC$,
contradicting irreducibility. If $\beta \neq 0$, the second intersection is $s = -\alpha/\beta \in \QQ$, a rational point (equal to $P$ when $\alpha = 0$, the tangent line). If $\beta = 0$, the line meets $C$
only at $P$; this happens for the directions $(u, v)$ with $f_2(u, v) = 0$, where $f_2$ is the degree-$2$ part of $f$, and there are at most two of them up to scaling. Conversely, a rational point $Q \neq P$ of $C$
lies on the rational line $PQ$, and on no other line through $P$. Since there are infinitely many rational directions and at most three of them (two asymptotic, one tangent) give no point other than $P$,
$C(\QQ)$ is infinite.
\end{proof}

So a conic either has no rational point or has as many as the projective line. Which conics have a rational point is decided by the Hasse--Minkowski theorem
(Theorem~\ref{p-adic-numbers:theorem-hasse-minkowski}): after completing squares one is reduced to a diagonal form $aX^2 + bY^2 + cZ^2 = 0$, which has a nontrivial rational solution if and only if it has one in
$\RR$ and in every $\QQ_p$.

\begin{example}\label{diophantine-equations:example-conics}
\begin{enumerate}
\item $X^2 + Y^2 = 3$ has no rational point. If $(a/c)^2 + (b/c)^2 = 3$ with $\gcd(a, b, c) = 1$, then $a^2 + b^2 = 3c^2$; squares are $0$ or $1$ modulo $3$, so $3 \mid a$ and $3 \mid b$, then $9 \mid 3c^2$ and
$3 \mid c$, a contradiction. This is a $3$-adic obstruction: the form $X^2 + Y^2 - 3Z^2$ is anisotropic over $\QQ_3$.
\item $X^2 + Y^2 = 2$ has the point $(1, 1)$, hence infinitely many: lines through $(1, 1)$ with slope $t$ give $\bigl(\frac{t^2 - 2t - 1}{t^2 + 1}, \frac{-t^2 - 2t + 1}{t^2 + 1}\bigr)$, for instance
$(\frac15, -\frac75)$ at $t = 3$.
\item $X^2 - 2Y^2 = 1$ has the rational points obtained from $(1, 0)$; its \emph{integer} points are the solutions of Pell's equation, which form an infinite cyclic group up to sign
(Theorem~\ref{continued-fractions:theorem-pell}). Integer points on conics are a much subtler question than rational points.
\end{enumerate}
\end{example}

\section{Infinite descent}\label{diophantine-equations:section-descent}

Fermat's method of \emph{infinite descent} proves that an equation has no solutions in positive integers by producing, from any solution, another one which is smaller in some positive integer measure; since
positive integers cannot decrease forever, there is no solution. Equivalently, one takes a solution with the measure minimal and derives a contradiction.

\begin{theorem}[Fermat]\label{diophantine-equations:theorem-four-plus}
The equation $x^4 + y^4 = z^2$ has no solution in positive integers.
\end{theorem}

\begin{proof}
Suppose there is one, and take one with $z$ minimal.

\emph{Step 1: reduction to a primitive triple.} If a prime $p$ divides $x$ and $y$, then $p^4 \mid z^2$, so $p^2 \mid z$ and $(x/p, y/p, z/p^2)$ is a smaller solution. Hence $\gcd(x, y) = 1$, and
$(x^2, y^2, z)$ is a primitive Pythagorean triple. Exchanging $x$ and $y$, assume $y$ even. By Theorem~\ref{diophantine-equations:theorem-pythagorean},
\[
x^2 = m^2 - n^2, \qquad y^2 = 2mn, \qquad z = m^2 + n^2,
\]
with $m > n > 0$ coprime of opposite parity.

\emph{Step 2: a second triple.} From $x^2 + n^2 = m^2$ with $x$ odd: if $m$ were even and $n$ odd, then $x^2 = m^2 - n^2 \equiv -1 \pmod 4$, impossible. So $n$ is even, and $(x, n, m)$ is a primitive triple
with $n$ even: $x = r^2 - s^2$, $n = 2rs$, $m = r^2 + s^2$ with $r > s > 0$ coprime.

\emph{Step 3: three squares.} Now $(y/2)^2 = \frac{2mn}{4} = mrs$, and $m$, $r$, $s$ are pairwise coprime ($r$, $s$ are coprime, and a prime dividing $r$ and $m = r^2 + s^2$ divides $s$). By
Lemma~\ref{diophantine-equations:lemma-coprime-powers}, $m = c^2$, $r = a^2$, $s = b^2$ with positive integers $a$, $b$, $c$.

\emph{Step 4: descent.} Then $a^4 + b^4 = r^2 + s^2 = m = c^2$, a new solution, with $c \leq c^2 = m \leq m^2 < m^2 + n^2 = z$. This contradicts the minimality of $z$.
\end{proof}

\begin{corollary}\label{diophantine-equations:corollary-flt-four}
For every $n$ divisible by $4$, the equation $x^n + y^n = z^n$ has no solution in positive integers.
\end{corollary}

\begin{proof}
If $n = 4k$, a solution gives $(x^k)^4 + (y^k)^4 = (z^{2k})^2$, contradicting Theorem~\ref{diophantine-equations:theorem-four-plus}.
\end{proof}

\begin{theorem}\label{diophantine-equations:theorem-four-minus}
The equation $x^4 - y^4 = z^2$ has no solution in positive integers.
\end{theorem}

\begin{proof}
Take a solution with $x$ minimal. As in Step 1 above, $\gcd(x, y) = 1$, and $z^2 + (y^2)^2 = (x^2)^2$ is a primitive Pythagorean triple; in particular $x$ is odd.

\emph{Case $y$ odd.} Then $z$ is the even leg: $y^2 = m^2 - n^2$, $z = 2mn$, $x^2 = m^2 + n^2$ with $m > n > 0$ coprime. Hence $m^4 - n^4 = (m^2 - n^2)(m^2 + n^2) = (xy)^2$, a solution with
$m < \sqrt{m^2 + n^2} = x$.

\emph{Case $y$ even.} Then $y^2 = 2mn$, $z = m^2 - n^2$, $x^2 = m^2 + n^2$ with $m > n > 0$ coprime of opposite parity. Now $(m, n, x)$ is a primitive triple, so $\{m, n\} = \{r^2 - s^2, 2rs\}$ and
$x = r^2 + s^2$ with $r > s > 0$ coprime of opposite parity. Then $(y/2)^2 = \frac{mn}{2} = rs(r - s)(r + s)$. The four factors are positive and pairwise coprime: $r$, $s$ are coprime; $r \pm s$ are odd and
$\gcd(r - s, r + s)$ divides $2r$ and $2s$, hence is $1$; and a prime dividing $r$ and $r \pm s$ divides $s$. By Lemma~\ref{diophantine-equations:lemma-coprime-powers}, $r = a^2$, $s = b^2$, $r + s = c^2$,
$r - s = d^2$, and $a^4 - b^4 = r^2 - s^2 = (cd)^2$, a solution with $a \leq a^2 = r < r^2 + s^2 = x$.

In both cases $x$ was not minimal, a contradiction.
\end{proof}

\begin{definition}\label{diophantine-equations:definition-congruent}
A positive rational number $n$ is a \emph{congruent number} if it is the area of a right triangle with rational sides, that is, if there are $a, b, c \in \QQ_{> 0}$ with $a^2 + b^2 = c^2$ and $\frac12ab = n$.
\end{definition}

Scaling the triangle by $\lambda \in \QQ^\times$ multiplies the area by $\lambda^2$, so whether $n$ is congruent depends only on its class modulo squares, and one may restrict to squarefree integers $n$.

\begin{example}\label{diophantine-equations:example-congruent}
$6$ is congruent ($3, 4, 5$), and so is $5$, with sides $\frac32, \frac{20}{3}, \frac{41}{6}$: indeed $\frac94 + \frac{400}{9} = \frac{1681}{36} = (\frac{41}{6})^2$ and $\frac12 \cdot \frac32 \cdot \frac{20}{3} = 5$.
$7$ is congruent, with sides $\frac{24}{5}, \frac{35}{12}, \frac{337}{60}$. Every primitive triple $(m^2 - n^2, 2mn, m^2 + n^2)$ shows that $mn(m^2 - n^2)$ is congruent; $(2, 1)$ gives $6$, and $(5, 4)$
gives $180 = 5 \cdot 6^2$, which is how $5$ is found.
\end{example}

\begin{theorem}[Fermat's right triangle theorem]\label{diophantine-equations:theorem-right-triangle}
No right triangle with rational sides has an area which is the square of a rational number. Equivalently, $1$ is not a congruent number.
\end{theorem}

\begin{proof}
\emph{Step 1: reduction.} Suppose a rational right triangle has square area. Scaling by a common denominator and then dividing by the greatest common divisor of the sides (both change the area by a square
factor) gives a primitive Pythagorean triple whose area is a square.

\emph{Step 2: four squares.} By Theorem~\ref{diophantine-equations:theorem-pythagorean} the legs are $m^2 - n^2$ and $2mn$, and the area is $mn(m - n)(m + n)$, with $m > n > 0$ coprime of opposite parity. As in
the proof of Theorem~\ref{diophantine-equations:theorem-four-minus}, the four factors are positive and pairwise coprime, so they are squares: $m = p^2$, $n = q^2$, $m + n = u^2$, $m - n = v^2$.

\emph{Step 3: conclusion.} Then $p^4 - q^4 = m^2 - n^2 = (uv)^2$ with $p, q, uv > 0$, contradicting Theorem~\ref{diophantine-equations:theorem-four-minus}.
\end{proof}

Congruent numbers are the first place where elliptic curves enter elementary number theory.

\begin{proposition}\label{diophantine-equations:proposition-congruent-curve}
A positive rational $n$ is congruent if and only if the curve $E_n : y^2 = x^3 - n^2x$ has a rational point with $y \neq 0$. The maps
\[
(a, b, c) \mapsto \Bigl(\frac{c^2}{4}, \frac{c(a^2 - b^2)}{8}\Bigr), \qquad (x, y) \mapsto \Bigl(\Bigl|\frac{x^2 - n^2}{y}\Bigr|, \Bigl|\frac{2nx}{y}\Bigr|, \Bigl|\frac{x^2 + n^2}{y}\Bigr|\Bigr)
\]
go between triangles and such points.
\end{proposition}

\begin{proof}
\emph{From triangles to points.} Let $a^2 + b^2 = c^2$ and $ab = 2n$, and put $x = c^2/4$. Then $x \pm n = \frac{c^2 \pm 2ab}{4} = \bigl(\frac{a \pm b}{2}\bigr)^2$, so
$x^3 - n^2x = x(x - n)(x + n) = \bigl(\frac{c(a - b)(a + b)}{8}\bigr)^2$. The value $y = \frac{c(a^2 - b^2)}{8}$ is nonzero, since $a = b$ would give $c = a\sqrt2$, which is irrational.

\emph{From points to triangles.} Let $y^2 = x^3 - n^2x$ with $y \neq 0$, and $a = \frac{x^2 - n^2}{y}$, $b = \frac{2nx}{y}$, $c = \frac{x^2 + n^2}{y}$. Then
$a^2 + b^2 = \frac{(x^2 - n^2)^2 + 4n^2x^2}{y^2} = \frac{(x^2 + n^2)^2}{y^2} = c^2$ and $\frac12ab = \frac{nx(x^2 - n^2)}{y^2} = n\frac{x^3 - n^2x}{y^2} = n$. The sides are nonzero ($x^2 - n^2 \neq 0$ and
$x \neq 0$, as $y \neq 0$), and replacing them by their absolute values gives a triangle of area $n$.
\end{proof}

\begin{corollary}\label{diophantine-equations:corollary-y2-x3-x}
The only rational points of $y^2 = x^3 - x$ are $(0, 0)$ and $(\pm1, 0)$.
\end{corollary}

\begin{proof}
A rational point with $y \neq 0$ would make $1$ congruent, against Theorem~\ref{diophantine-equations:theorem-right-triangle}; the points with $y = 0$ are the roots of $x^3 - x$.
\end{proof}

\begin{remark}\label{diophantine-equations:remark-tunnell}
Deciding which $n$ are congruent is an open problem in general. The rational points of $E_n$ form a finitely generated abelian group (the Mordell--Weil theorem), and $n$ is congruent if and only if this group
is infinite. Tunnell gave in 1983 a simple criterion, counting representations of $n$ by certain ternary quadratic forms, which is necessary for $n$ to be congruent and is sufficient if the Birch and
Swinnerton-Dyer conjecture holds for $E_n$ \cite[Chapter IV]{koblitz-elliptic}. These topics belong to Chapter~\ref{elliptic-curves-arithmetic}.
\end{remark}

\section{Sums of three squares}\label{diophantine-equations:section-three-squares}

Lagrange's theorem (Theorem~\ref{number-theory:theorem-four-squares}) writes every positive integer as a sum of four squares, and Chapter~\ref{number-theory} described the sums of two squares. Sums of three
squares are harder. Chapter~\ref{p-adic-numbers} decided, by the Hasse--Minkowski theorem, which integers are sums of three \emph{rational} squares
(Theorem~\ref{p-adic-numbers:theorem-three-squares-rational}); a lemma of Davenport and Cassels turns rational representations into integer ones.

\begin{lemma}[Davenport--Cassels]\label{diophantine-equations:lemma-davenport-cassels}
Let $f(x) = \sum_{i, j = 1}^{n} a_{ij}x_ix_j$ be a positive definite quadratic form with $a_{ij} = a_{ji} \in \ZZ$, and write $B(x, y) = \sum a_{ij}x_iy_j$, so that $f(x) = B(x, x)$. Assume that for every
$x \in \QQ^n$ there is $y \in \ZZ^n$ with $f(x - y) < 1$. If an integer $m$ is represented by $f$ over $\QQ$, it is represented by $f$ over $\ZZ$.
\end{lemma}

\begin{proof}
\emph{Step 1: a minimal denominator.} Among the $x \in \QQ^n$ with $f(x) = m$ choose one whose common denominator is smallest: $x = z/t$ with $z \in \ZZ^n$ and $t \geq 1$ minimal. We show $t = 1$. Suppose
$t > 1$, choose $y \in \ZZ^n$ with $f(x - y) < 1$ and put $w = x - y$. Then $w \neq 0$, since $x \notin \ZZ^n$, so $0 < f(w) < 1$ by positive definiteness.

\emph{Step 2: a new denominator.} Put $a = f(y) - m \in \ZZ$ and $b = 2(mt - B(z, y)) \in \ZZ$. Then
\[
tf(w) = t\bigl(f(x) - 2B(x, y) + f(y)\bigr) = tm - 2B(z, y) + tf(y) = at + b.
\]
So $t^{\prime} = at + b = tf(w)$ is an integer with $0 < t^{\prime} < t$.

\emph{Step 3: a smaller denominator.} The point $x^{\prime} = y + \frac{a}{f(w)}w$ satisfies $f(x^{\prime}) = m$: the polynomial $\lambda \mapsto f(y + \lambda w) - m = f(w)\lambda^2 + 2B(y, w)\lambda + a$ has
the root $\lambda = 1$, and the product of its roots is $a/f(w)$, so $\lambda = a/f(w)$ is also a root. And
\[
x^{\prime} = y + \frac{at}{t^{\prime}}\Bigl(\frac zt - y\Bigr) = \frac{az + (t^{\prime} - at)y}{t^{\prime}} = \frac{az + by}{t^{\prime}},
\]
which has denominator dividing $t^{\prime} < t$. This contradicts the minimality of $t$.
\end{proof}

\begin{theorem}[Legendre, Gauss]\label{diophantine-equations:theorem-three-squares}
A positive integer $n$ is a sum of three squares of integers if and only if it is not of the form $4^a(8b + 7)$ with $a, b \geq 0$.
\end{theorem}

\begin{proof}
By Theorem~\ref{p-adic-numbers:theorem-three-squares-rational}, $n$ is a sum of three rational squares if and only if it is not of the form $4^a(8b + 7)$. The form $f = x_1^2 + x_2^2 + x_3^2$ is positive definite
with integer coefficients, and for $x \in \QQ^3$, rounding each coordinate to a nearest integer gives $y \in \ZZ^3$ with $f(x - y) \leq 3 \cdot \frac14 < 1$. By
Lemma~\ref{diophantine-equations:lemma-davenport-cassels}, rational and integer representations coincide.
\end{proof}

\begin{corollary}[Gauss]\label{diophantine-equations:corollary-triangular}
Every nonnegative integer is a sum of three triangular numbers $\frac{k(k + 1)}{2}$, $k \geq 0$.
\end{corollary}

\begin{proof}
$8n + 3$ is not of the form $4^a(8b + 7)$ (it is odd and $\equiv 3 \pmod 8$), so $8n + 3 = x^2 + y^2 + z^2$. Squares are $0$, $1$ or $4$ modulo $8$, and a sum of three of them is $\equiv 3 \pmod 8$ only as
$1 + 1 + 1$; so $x$, $y$, $z$ are odd, $x = 2k + 1$ and so on with $k, l, m \geq 0$ (replace $x$ by $-x$ if needed). Then $8n + 3 = 4k(k + 1) + 4l(l + 1) + 4m(m + 1) + 3$, and dividing by $8$ gives
$n = \frac{k(k + 1)}{2} + \frac{l(l + 1)}{2} + \frac{m(m + 1)}{2}$.
\end{proof}

\begin{example}\label{diophantine-equations:example-three-squares}
$7$, $15$, $23$, $28 = 4 \cdot 7$, $31$ are not sums of three squares, while $6 = 4 + 1 + 1$, $14 = 9 + 4 + 1$, $29 = 25 + 4$. Gauss recorded the corollary in his diary as "EYPHKA! num $= \Delta + \Delta + \Delta$";
for instance $10 = 6 + 3 + 1$ and $20 = 10 + 10 + 0$. The problem of writing $n$ as a sum of a bounded number of $k$-th powers is Waring's problem
(Example~\ref{additive-number-theory:example-waring}).
\end{example}

\begin{counterexample}\label{diophantine-equations:counterexample-davenport-cassels}
The rounding hypothesis of Lemma~\ref{diophantine-equations:lemma-davenport-cassels} is needed. For $f = x_1^2 + 14x_2^2$,
\[
2 = \Bigl(\frac23\Bigr)^2 + 14\Bigl(\frac13\Bigr)^2,
\]
since $4 + 14 = 18 = 2 \cdot 9$; but $2 = x_1^2 + 14x_2^2$ has no integer solution. Here $x = (\frac23, \frac13)$ has $f(x - y) \geq \frac19 + \frac{14}{9} > 1$ for every $y \in \ZZ^2$.
\end{counterexample}

\section{Markov's equation}\label{diophantine-equations:section-markov}

The equation
\begin{equation}\label{diophantine-equations:equation-markov}
x^2 + y^2 + z^2 = 3xyz
\end{equation}
arose in Markov's work on the worst approximable quadratic irrationals. Its positive solutions are \emph{Markov triples}, and their entries \emph{Markov numbers}. The key observation is \emph{Vieta jumping}:
for fixed $x$, $y$, the equation is quadratic in $z$, $z^2 - 3xyz + x^2 + y^2 = 0$, so with a root $z$ comes the root
\[
z^{\prime} = 3xy - z = \frac{x^2 + y^2}{z},
\]
which is again a positive integer. The \emph{move} $(x, y, z) \mapsto (x, y, 3xy - z)$, and its analogues in the other coordinates, send Markov triples to Markov triples.

\begin{theorem}[Markov]\label{diophantine-equations:theorem-markov}
Every Markov triple is obtained from $(1, 1, 1)$ by a finite sequence of moves and permutations of the coordinates.
\end{theorem}

\begin{proof}
\emph{Step 1: triples with repeated entries.} Let $(x, y, z)$ be a Markov triple, ordered $x \leq y \leq z$. If $y = z$, then $x^2 + 2y^2 = 3xy^2$, so $y^2 \mid x^2$, hence $y \leq x$, so $x = y = z$ and
$3x^2 = 3x^3$, $x = 1$. If $x = y$, then $2x^2 + z^2 = 3x^2z$, so $x^2 \mid z^2$, $z = kx$, and $2 + k^2 = 3kx$, so $k \mid 2$: $k = 1$ gives $(1, 1, 1)$ and $k = 2$ gives $(1, 1, 2)$.

\emph{Step 2: descent.} Suppose $x < y < z$. Let $g(t) = t^2 - 3xyt + x^2 + y^2$, whose roots are $z$ and $z^{\prime} = 3xy - z$. Then
$g(y) = 2y^2 + x^2 - 3xy^2 < 3y^2 - 3xy^2 \leq 0$, as $x < y$ and $x \geq 1$. So $y$ lies strictly between the roots: $z^{\prime} < y < z$. The move in the last coordinate gives the triple $(x, y, z^{\prime})$,
whose largest entry $y$ is smaller than $z$.

\emph{Step 3: conclusion.} Repeating Step 2, the largest entry decreases strictly, so after finitely many moves and reorderings one reaches a triple with a repeated entry, which by Step 1 is $(1, 1, 1)$ or
$(1, 1, 2)$; and $(1, 1, 2)$ is the move of $(1, 1, 1)$ in the last coordinate. Reversing the moves (each move is an involution) gives the claim.
\end{proof}

\begin{example}\label{diophantine-equations:example-markov}
Starting from $(1, 1, 2)$, each triple $(x, y, z)$ with $z$ largest has two "children" obtained by moving $x$ or $y$, which gives a binary tree, the \emph{Markov tree}:
\[
\begin{aligned}
&(1, 1, 2) \to (1, 2, 5) \to (1, 5, 13),\ (2, 5, 29), \\
&(1, 5, 13) \to (1, 13, 34),\ (5, 13, 194), \qquad (2, 5, 29) \to (2, 29, 169),\ (5, 29, 433).
\end{aligned}
\]
The Markov numbers begin $1, 2, 5, 13, 29, 34, 89, 169, 194, 233, 433, \dots$; the branch $(1, F_{2k - 1}, F_{2k + 1})$ consists of odd-indexed Fibonacci numbers, and the branch $(2, P_{2k - 1}, P_{2k + 1})$ of
odd-indexed Pell numbers.
\end{example}

\begin{remark}\label{diophantine-equations:remark-markov}
Markov proved that the irrationals $\alpha$ for which $|\alpha - p/q| < \frac{1}{c\,q^2}$ has infinitely many solutions for every $c$ smaller than a given $c_\alpha$ exhaust, below $3$, the values
$c = \sqrt{9m^2 - 4}/m$ for $m$ a Markov number: $\sqrt5$, $\sqrt8$, $\sqrt{221}/5$, \dots, accumulating at $3$; the first is Hurwitz's constant from Chapter~\ref{continued-fractions}. Frobenius conjectured
in 1913 that a Markov number is the largest entry of only one Markov triple up to order; this \emph{uniqueness conjecture} is still open \cite{aigner-markov}.
\end{remark}

\section{Mordell's equation and other classical equations}\label{diophantine-equations:section-mordell}

\emph{Mordell's equation} is $y^2 = x^3 + k$ with $k \in \ZZ \setminus \{0\}$. For each $k$ it has only finitely many integer solutions \cite{mordell-diophantine}, but there is no uniform elementary method;
two typical cases follow, one solved by congruences and one by factorisation in a ring of quadratic integers.

\begin{proposition}[Lebesgue]\label{diophantine-equations:proposition-mordell-seven}
The equation $y^2 = x^3 + 7$ has no integer solution.
\end{proposition}

\begin{proof}
\emph{$x$ is odd.} If $x$ is even, $y^2 \equiv 7 \pmod 8$, while odd squares are $\equiv 1 \pmod 8$.

\emph{A prime $q \equiv 3 \pmod 4$.} Write $y^2 + 1 = x^3 + 8 = (x + 2)(x^2 - 2x + 4)$. As $x$ is odd, $x^2 - 2x + 4 = (x - 1)^2 + 3 \equiv 3 \pmod 4$, and it is positive, so it has a prime factor
$q \equiv 3 \pmod 4$. Then $y^2 \equiv -1 \pmod q$, which is impossible by Corollary~\ref{number-theory:corollary-minus-one}.
\end{proof}

\begin{proposition}[Fermat, Euler]\label{diophantine-equations:proposition-mordell-two}
The only integer solutions of $y^2 = x^3 - 2$ are $(x, y) = (3, \pm5)$.
\end{proposition}

\begin{proof}
\emph{Step 1: the ring $\ZZ[\sqrt{-2}]$.} It is Euclidean for the norm $N(a + b\sqrt{-2}) = a^2 + 2b^2$: given $\alpha, \beta \neq 0$, write $\alpha/\beta = u + v\sqrt{-2}$ with $u, v \in \QQ$, choose integers $m$,
$n$ with $|u - m|, |v - n| \leq \frac12$, and put $\kappa = m + n\sqrt{-2}$; then $N(\alpha/\beta - \kappa) \leq \frac14 + \frac24 < 1$, so $N(\alpha - \kappa\beta) < N(\beta)$. Hence it is a UFD
(Theorem~\ref{rings:theorem-euclidean-pid-ufd}), and its units, the elements of norm $1$, are $\pm1$, both cubes.

\emph{Step 2: coprime factors.} Let $y^2 + 2 = x^3$. Then $y$ is odd: if $y$ were even, $x$ would be even, and $8 \mid x^3$ while $y^2 + 2 \equiv 2 \pmod 4$. Factor
$x^3 = (y + \sqrt{-2})(y - \sqrt{-2})$. A prime $\delta$ dividing both factors divides their difference $2\sqrt{-2} = -(\sqrt{-2})^3$, so $\delta$ is associate to the prime $\sqrt{-2}$ (of norm $2$), and then
$2 = N(\sqrt{-2})$ divides $N(y + \sqrt{-2}) = y^2 + 2$, which is odd. So the factors are coprime.

\emph{Step 3: a cube.} By Lemma~\ref{diophantine-equations:lemma-coprime-powers}, $y + \sqrt{-2}$ is a unit times a cube, hence a cube, as $\pm1 = (\pm1)^3$: $y + \sqrt{-2} = (a + b\sqrt{-2})^3$ with
$a, b \in \ZZ$, so
\[
y + \sqrt{-2} = (a^3 - 6ab^2) + (3a^2b - 2b^3)\sqrt{-2}.
\]
Comparing coefficients of $\sqrt{-2}$: $b(3a^2 - 2b^2) = 1$, so $b = \pm1$ and $3a^2 - 2 = b$; this forces $b = 1$, $a = \pm1$. Then $y = a^3 - 6a = \mp5$ and $x^3 = 25 + 2 = 27$, $x = 3$.
\end{proof}

\begin{remark}\label{diophantine-equations:remark-mordell}
The same method applied to $y^2 = x^3 - k$ needs unique factorisation in $\ZZ[\sqrt{-k}]$, which often fails (for $k = 5$, Remark~\ref{number-systems:remark-unique-factorisation-fails}); the repair, by ideals
and the finiteness of the class number, is the subject of Chapter~\ref{algebraic-number-theory}. Equations of the form $F(x, y) = m$ with $F$ an irreducible binary form of degree at least $3$ (\emph{Thue
equations}, such as $x^3 - 2y^3 = 1$) have only finitely many integer solutions; Thue deduced this in 1909 from his improvement of Liouville's approximation theorem, and Baker made it effective
\cite{baker-transcendental}. Mordell's equation reduces to finitely many Thue equations, which gives its finiteness.
\end{remark}

\begin{theorem}[Mih\u{a}ilescu]\label{diophantine-equations:theorem-catalan}
The only solution of $x^m - y^n = 1$ in integers $x, y, m, n \geq 2$ is $3^2 - 2^3 = 1$ \cite{mihailescu-2004}.
\end{theorem}

Catalan conjectured this in 1844. The special case of powers of $2$ and $3$ is elementary (Exercise~\ref{diophantine-equations:exercise-powers-two-three}).

\section{Fermat's last theorem}\label{diophantine-equations:section-flt}

\begin{theorem}[Wiles, Taylor--Wiles]\label{diophantine-equations:theorem-flt}
For $n \geq 3$, the equation $x^n + y^n = z^n$ has no solution in positive integers \cite{wiles-1995}, \cite{taylor-wiles-1995}.
\end{theorem}

Fermat stated this in a margin around 1637; the proof came in 1994, through the modularity of semistable elliptic curves (Chapter~\ref{elliptic-curves-arithmetic} and Chapter~\ref{modular-forms}). Since
every $n \geq 3$ is divisible by $4$ or by an odd prime, and a solution for $n = dk$ gives one for $d$, it is enough to treat $n = 4$, done in Corollary~\ref{diophantine-equations:corollary-flt-four}, and odd
primes $n = p$. This section proves the case $p = 3$ and Sophie Germain's theorem, and the analogue for polynomials.

\subsection{The exponent three}\label{diophantine-equations:subsection-three}

We work in the Eisenstein integers $\ZZ[\omega]$, $\omega = e^{2\pi i/3}$, a UFD with units $\pm1, \pm\omega, \pm\omega^2$ (Proposition~\ref{number-theory:proposition-eisenstein}). Let $\lambda = 1 - \omega$, a
prime of norm $3$ with $3 = -\omega^2\lambda^2$.

\begin{lemma}\label{diophantine-equations:lemma-eisenstein-cubes}
\begin{enumerate}
\item Every element of $\ZZ[\omega]$ is congruent to $0$, $1$ or $-1$ modulo $\lambda$.
\item If $\alpha \equiv \pm1 \pmod \lambda$, then $\alpha^3 \equiv \pm1 \pmod{\lambda^4}$, with the same sign.
\item A rational integer is divisible by $\lambda$ if and only if it is divisible by $3$.
\end{enumerate}
\end{lemma}

\begin{proof}
(1) $\omega \equiv 1 \pmod \lambda$, so $a + b\omega \equiv a + b \pmod \lambda$; and $\lambda \mid 3$, so an integer is congruent modulo $\lambda$ to its remainder modulo $3$.

(3) If $\lambda \mid k \in \ZZ$, then $3 = N(\lambda)$ divides $N(k) = k^2$, so $3 \mid k$; the converse holds as $\lambda \mid 3$.

(2) It suffices to treat $\alpha = 1 + \beta\lambda$ (replace $\alpha$ by $-\alpha$). Then
$\alpha^3 - 1 = (\alpha - 1)(\alpha - \omega)(\alpha - \omega^2)$, and $\alpha - \omega = \beta\lambda + \lambda$, $\alpha - \omega^2 = \beta\lambda + (1 - \omega^2) = \beta\lambda - \omega^2\lambda$, using
$1 - \omega^2 = (1 - \omega)(1 + \omega) = -\omega^2\lambda$. So $\alpha^3 - 1 = \lambda^3\beta(\beta + 1)(\beta - \omega^2)$. As $\omega^2 \equiv 1 \pmod \lambda$, the three factors are congruent to
$\beta$, $\beta + 1$, $\beta - 1$, and by (1) one of them is divisible by $\lambda$. Hence $\lambda^4 \mid \alpha^3 - 1$.
\end{proof}

\begin{theorem}[Euler, Gauss]\label{diophantine-equations:theorem-flt-three}
The equation $x^3 + y^3 = z^3$ has no solution in nonzero integers.
\end{theorem}

\begin{proof}
Suppose $x^3 + y^3 = z^3$ with $xyz \neq 0$; dividing by common factors, $x$, $y$, $z$ are pairwise coprime. Writing the equation as $x^3 + y^3 + (-z)^3 = 0$, it is symmetric in the three unknowns up to sign.

\emph{Step 1: $3 \nmid xyz$ is impossible.} Cubes of integers prime to $3$ are $\pm1 \pmod 9$ (check $1, 2, 4, 5, 7, 8$), so $x^3 + y^3 + (-z)^3 \equiv \pm1 \pm 1 \pm 1 \not\equiv 0 \pmod 9$.

\emph{Step 2: reformulation.} So $3$ divides exactly one of the unknowns; by the symmetry, say $3 \mid z$, $z = 3^kw$ with $k \geq 1$, $3 \nmid w$. Then $z^3 = (-\omega^2\lambda^2)^{3k}w^3 = \varepsilon\lambda^{6k}w^3$
with $\varepsilon$ a unit. By Lemma~\ref{diophantine-equations:lemma-eisenstein-cubes}(3), $\lambda \nmid xyw$, and $x$, $y$ are coprime in $\ZZ[\omega]$ (a Bézout relation $ux + vy = 1$ in $\ZZ$ persists). So it
suffices to prove:

\emph{Claim.} There are no $\alpha, \beta, \gamma \in \ZZ[\omega]$, unit $\varepsilon$ and integer $n \geq 1$ with $\alpha^3 + \beta^3 = \varepsilon\lambda^{3n}\gamma^3$, $\alpha$, $\beta$ coprime, and
$\lambda \nmid \alpha\beta\gamma$.

Take a counterexample with $n$ minimal.

\emph{Step 3: $n \geq 2$.} By Lemma~\ref{diophantine-equations:lemma-eisenstein-cubes}, $\alpha^3 \equiv a$ and $\beta^3 \equiv b \pmod{\lambda^4}$ with $a, b \in \{\pm1\}$. Since $\lambda \mid \alpha^3 + \beta^3$,
$\lambda \mid a + b \in \{0, \pm2\}$; as $\lambda \nmid 2$ (its norm $3$ does not divide $4$), $a + b = 0$, so $\lambda^4 \mid \alpha^3 + \beta^3$ and $3n \geq 4$.

\emph{Step 4: factorisation.} $\alpha^3 + \beta^3 = (\alpha + \beta)(\alpha + \omega\beta)(\alpha + \omega^2\beta)$. The differences of the factors are $\lambda\beta$ times units, and
$\omega(\alpha + \beta) - (\alpha + \omega\beta) = -\lambda\alpha$; so each factor is divisible by $\lambda$ (one is, since $\lambda$ divides the product, and then all are), and a common prime factor of two
of them divides $\lambda\alpha$ and $\lambda\beta$, hence is $\lambda$, and $\lambda^2$ does not divide both (else $\lambda^2 \mid \lambda\beta$). So two factors are exactly divisible by $\lambda$ and the third by
$\lambda^{3n - 2}$. Replacing $\beta$ by $\omega\beta$ or $\omega^2\beta$, which changes nothing in the hypotheses, assume the third factor is $\alpha + \beta$. The three quotients
$\frac{\alpha + \beta}{\lambda^{3n - 2}}$, $\frac{\alpha + \omega\beta}{\lambda}$, $\frac{\alpha + \omega^2\beta}{\lambda}$ are pairwise coprime, prime to $\lambda$, and their product is a unit times $\gamma^3$, so
by Lemma~\ref{diophantine-equations:lemma-coprime-powers}
\[
\alpha + \beta = \varepsilon_1\lambda^{3n - 2}\theta^3, \quad \alpha + \omega\beta = \varepsilon_2\lambda\psi^3, \quad \alpha + \omega^2\beta = \varepsilon_3\lambda\varphi^3,
\]
with units $\varepsilon_i$ and $\theta$, $\psi$, $\varphi$ pairwise coprime and prime to $\lambda$.

\emph{Step 5: a smaller solution.} Since $1 + \omega + \omega^2 = 0$, $(\alpha + \beta) + \omega(\alpha + \omega\beta) + \omega^2(\alpha + \omega^2\beta) = 0$. Dividing by $\omega\varepsilon_2\lambda$:
\[
\psi^3 + \varepsilon_4\varphi^3 = \varepsilon_5\lambda^{3(n - 1)}\theta^3
\]
for units $\varepsilon_4$, $\varepsilon_5$. As $n \geq 2$, $\lambda^3$ divides the left side, and $\psi^3, \varphi^3 \equiv \pm1 \pmod{\lambda^4}$, so $\pm1 \pm \varepsilon_4 \equiv 0 \pmod{\lambda^2}$. Among the
units only $\varepsilon_4 = \pm1$ allow this: $1 \pm \omega$ and $1 \pm \omega^2$ are units or $\lambda$ times units. Replacing $\varphi$ by $-\varphi$ if $\varepsilon_4 = -1$, we get
$\psi^3 + \varphi^3 = \varepsilon_5\lambda^{3(n - 1)}\theta^3$, a counterexample with $n - 1 \geq 1$, contradicting minimality.
\end{proof}

\subsection{Sophie Germain's theorem}\label{diophantine-equations:subsection-germain}

For a general odd prime $p$ one distinguishes the \emph{first case}, $p \nmid xyz$, from the second. Sophie Germain found a simple criterion for the first case.

\begin{theorem}[Sophie Germain]\label{diophantine-equations:theorem-germain}
Let $p$ be an odd prime such that $q = 2p + 1$ is also prime. Then $x^p + y^p = z^p$ has no solution in integers with $p \nmid xyz$.
\end{theorem}

\begin{proof}
Suppose there is one. As $p$ is odd, we may write it as $x^p + y^p + z^p = 0$ (replace $z$ by $-z$), with $x$, $y$, $z$ pairwise coprime and $p \nmid xyz$.

\emph{Step 1: $p$-th roots.} We have $-x^p = (y + z)Q$ with $Q = \frac{y^p + z^p}{y + z} = \sum_{i=0}^{p-1}(-1)^iy^{p-1-i}z^i \equiv py^{p-1} \pmod{y + z}$. A prime $r$ dividing $y + z$ and $Q$ divides
$py^{p-1}$; it does not divide $y$ (it would divide $z$), so $r = p$, and then $p \mid y + z \mid x^p$, against $p \nmid x$. So $y + z$ and $Q$ are coprime, and since their product is a $p$-th power with $p$
odd, $y + z = a^p$ and $Q = \alpha^p$ for integers $a$, $\alpha$ (Lemma~\ref{diophantine-equations:lemma-coprime-powers}, signs absorbed into the odd powers). By symmetry also $z + x = b^p$ and $x + y = c^p$.

\emph{Step 2: $p$-th powers modulo $q$.} For $u \not\equiv 0 \pmod q$, $(u^p)^2 = u^{q - 1} \equiv 1$, so $u^p \equiv \pm1 \pmod q$. If $q$ divided none of $x$, $y$, $z$, then $x^p + y^p + z^p$ would be
$\equiv \pm1 \pm1 \pm1 \in \{\pm1, \pm3\} \pmod q$, which is nonzero as $q \geq 7$. So $q$ divides one of them, say $q \mid x$ (the argument is symmetric).

\emph{Step 3: $q \mid a$.} Now $2x = b^p + c^p - a^p \equiv 0 \pmod q$, so by Step 2 again one of $a$, $b$, $c$ is divisible by $q$. If $q \mid b$, then $q \mid z + x$, so $q \mid z$, contradicting
$\gcd(x, z) = 1$; similarly $q \nmid c$. So $q \mid a$, that is, $y \equiv -z \pmod q$.

\emph{Step 4: contradiction.} Then $\alpha^p = Q \equiv py^{p-1} \pmod q$. Also $y \equiv x + y = c^p \pmod q$, with $q \nmid y$ (else $q \mid x, y$), so $y \equiv \pm1$ and $y^{p - 1} \equiv 1 \pmod q$, as
$p - 1$ is even. Hence $\alpha^p \equiv p \pmod q$, and $q \nmid \alpha$, so $p \equiv \pm1 \pmod q$. But $0 < p - 1 < p + 1 < q$, so $q$ divides neither $p - 1$ nor $p + 1$, a contradiction.
\end{proof}

\begin{example}\label{diophantine-equations:example-germain}
The \emph{Sophie Germain primes} $p$, those with $2p + 1$ prime, begin $2, 3, 5, 11, 23, 29, 41, 53, 83, 89, \dots$; it is not known whether there are infinitely many. Germain's theorem settles the first case for
all of them, and the same argument with auxiliary primes of the form $q = 2kp + 1$ applies to many more exponents.
\end{example}

\begin{remark}\label{diophantine-equations:remark-kummer}
Kummer factored $x^p + y^p = \prod_{j=0}^{p-1}(x + \zeta^jy)$ in $\ZZ[\zeta]$, $\zeta = e^{2\pi i/p}$, as in Step 4 above. Unique factorisation fails in general (first for $p = 23$), but Kummer showed that the
argument goes through when $p$ does not divide the class number of $\QQ(\zeta)$ (\emph{regular primes}, such as all $p < 37$); this is developed in Chapter~\ref{algebraic-number-theory}.
\end{remark}

\subsection{Polynomials and the abc conjecture}\label{diophantine-equations:subsection-abc}

For polynomials there is a short and complete proof of Fermat's last theorem, through an inequality whose integer analogue is the $abc$ conjecture. Let $k$ be a field of characteristic $0$. For a nonzero
polynomial $f \in k[t]$, let $n_0(f)$ be the number of distinct roots of $f$ in an algebraic closure; equivalently, the degree of the product of the distinct monic irreducible factors of $f$, the \emph{radical}
$\operatorname{rad}(f)$. The derivative $f^{\prime}$ is the formal one.

\begin{theorem}[Mason--Stothers]\label{diophantine-equations:theorem-mason-stothers}
Let $a, b, c \in k[t]$ be pairwise coprime, not all constant, with $a + b = c$. Then
\[
\max(\deg a, \deg b, \deg c) \leq n_0(abc) - 1.
\]
\end{theorem}

\begin{proof}
\emph{Step 1: the Wronskian.} Let $W = ab^{\prime} - a^{\prime}b$. From $c = a + b$ and $c^{\prime} = a^{\prime} + b^{\prime}$ one gets $W = ac^{\prime} - a^{\prime}c = cb^{\prime} - c^{\prime}b$. Moreover $W \neq 0$:
if $ab^{\prime} = a^{\prime}b$, then $a \mid a^{\prime}b$, so $a \mid a^{\prime}$ by coprimality, so $a^{\prime} = 0$ by degree, and $a$ is constant since the characteristic is $0$; likewise $b$ is constant, and
so is $c$, a contradiction.

\emph{Step 2: divisibility.} If $\pi^m$ divides $a$ for an irreducible $\pi$, then $\pi^{m-1}$ divides $a^{\prime}$, so $\pi^{m - 1}$ divides $W = ab^{\prime} - a^{\prime}b$. Using the three expressions of $W$,
the same holds for the prime powers dividing $b$ and $c$. As $a$, $b$, $c$ are pairwise coprime, $\frac{abc}{\operatorname{rad}(abc)}$ divides $W$.

\emph{Step 3: degrees.} Hence $\deg a + \deg b + \deg c - n_0(abc) \leq \deg W \leq \deg a + \deg b - 1$, so $\deg c \leq n_0(abc) - 1$. The hypotheses are symmetric under permuting $a$, $b$, $-c$, which gives the
same bound for $\deg a$ and $\deg b$.
\end{proof}

\begin{corollary}\label{diophantine-equations:corollary-flt-polynomials}
For $n \geq 3$ there are no coprime polynomials $x, y, z \in k[t]$, not all constant, with $x^n + y^n = z^n$.
\end{corollary}

\begin{proof}
A common factor of two of them would divide the third, so they are pairwise coprime, and so are $x^n$, $y^n$, $z^n$. Let $d = \max(\deg x, \deg y, \deg z) \geq 1$. Since
$n_0(x^ny^nz^n) = n_0(xyz) \leq \deg x + \deg y + \deg z \leq 3d$, the theorem gives $nd \leq 3d - 1$, so $n < 3$.
\end{proof}

For integers, the radical $\operatorname{rad}(m)$ is the product of the distinct primes dividing $m$, and the analogy suggests the following.

\begin{conjecture}[$abc$ conjecture, Masser--Oesterl\'e]\label{diophantine-equations:conjecture-abc}
For every $\epsilon > 0$ there is a constant $K_\epsilon$ such that for all coprime positive integers $a$, $b$, $c$ with $a + b = c$,
\[
c \leq K_\epsilon\operatorname{rad}(abc)^{1 + \epsilon}.
\]
\end{conjecture}

\begin{remark}\label{diophantine-equations:remark-abc}
The exponent $1 + \epsilon$ cannot be replaced by $1$: for $a = 1$, $c = 3^{2^k}$, $b = c - 1$ with $k \geq 1$, induction on $k$ gives $2^{k + 2} \mid b$ (as $3^{2^{k+1}} - 1 = (3^{2^k} - 1)(3^{2^k} + 1)$), so $\operatorname{rad}(abc) \leq 3b/2^{k + 1} < 3c/2^{k + 1}$, and
$c/\operatorname{rad}(abc)$ is unbounded. The conjecture would imply that $x^n + y^n = z^n$ has no solutions for $n$ large (apply it to $x^n + y^n = z^n$: $z^n \leq K_\epsilon(xyz)^{1 + \epsilon} < K_\epsilon z^{3(1 + \epsilon)}$),
as well as many other finiteness results \cite{granville-tucker-abc}. A proof announced by Mochizuki in 2012 is not accepted by most experts.
\end{remark}

\section{Hilbert's tenth problem}\label{diophantine-equations:section-hilbert-tenth}

Hilbert's tenth problem (1900) asked for a procedure deciding, for any polynomial equation with integer coefficients, whether it has an integer solution. The answer is negative, and the reason is that
Diophantine equations are as expressive as computer programs.

\begin{definition}\label{diophantine-equations:definition-diophantine-set}
A set $S \subseteq \NN^k$ is \emph{Diophantine} if there is a polynomial $P \in \ZZ[a_1, \dots, a_k, x_1, \dots, x_m]$ such that
\[
(a_1, \dots, a_k) \in S \iff \exists x_1, \dots, x_m \in \NN : P(a_1, \dots, a_k, x_1, \dots, x_m) = 0.
\]
\end{definition}

\begin{example}\label{diophantine-equations:example-diophantine-sets}
\begin{enumerate}
\item The composite numbers: $a = (x_1 + 2)(x_2 + 2)$.
\item The relation $a_1 \leq a_2$: $a_2 = a_1 + x_1$. Divisibility $a_1 \mid a_2$: $a_2 = a_1x_1$.
\item Unknowns ranging over $\ZZ$ instead of $\NN$ give the same notion: $x \in \ZZ$ can be written $x_1 - x_2$ with $x_i \in \NN$, and conversely $x \in \NN$ if and only if $x = y_1^2 + y_2^2 + y_3^2 + y_4^2$
with $y_i \in \ZZ$ (Theorem~\ref{number-theory:theorem-four-squares}).
\end{enumerate}
\end{example}

\begin{proposition}\label{diophantine-equations:proposition-diophantine-closure}
Unions and intersections of two Diophantine subsets of $\NN^k$ are Diophantine.
\end{proposition}

\begin{proof}
If $S$ is defined by $P(a, x) = 0$ and $T$ by $Q(a, y) = 0$, with disjoint sets of unknowns $x$ and $y$, then $S \cup T$ is defined by $P(a, x)Q(a, y) = 0$ and $S \cap T$ by $P(a, x)^2 + Q(a, y)^2 = 0$.
\end{proof}

\begin{theorem}[Matiyasevich, Robinson, Davis, Putnam]\label{diophantine-equations:theorem-mrdp}
A subset of $\NN^k$ is Diophantine if and only if it is recursively enumerable, that is, it is the set of outputs of some algorithm \cite{matiyasevich-hilbert}.
\end{theorem}

\begin{corollary}\label{diophantine-equations:corollary-hilbert-tenth}
There is no algorithm which decides, for every polynomial $P \in \ZZ[x_1, \dots, x_m]$, whether $P = 0$ has a solution in integers.
\end{corollary}

\begin{proof}
There are recursively enumerable sets $S \subseteq \NN$ whose membership is undecidable (the set of programs that halt, suitably encoded) \cite{matiyasevich-hilbert}. By the theorem $S$ is defined by some
$P(a, x)$ with unknowns in $\NN$. Substituting for each $x_i$ a sum of four squares of new unknowns (Example~\ref{diophantine-equations:example-diophantine-sets}(3)) gives, for each $a$, a polynomial
equation with unknowns in $\ZZ$ which is solvable if and only if $a \in S$; an algorithm for Hilbert's problem would therefore decide membership in $S$.
\end{proof}

\begin{remark}\label{diophantine-equations:remark-hilbert-tenth}
Two striking consequences \cite{matiyasevich-hilbert}: there is a polynomial with integer coefficients whose positive values, as its variables range over $\NN$, are exactly the prime numbers (Jones, Sato, Wada
and Wiens wrote one of degree $25$ in $26$ variables); and there is a single polynomial $P(a, x_1, \dots, x_m)$, a "universal Diophantine equation", from which every Diophantine set is obtained by fixing one
parameter. The analogue of Hilbert's tenth problem over $\QQ$ is open. Computability itself is the subject of Chapter~\ref{computability-and-complexity}.
\end{remark}

\section{Exercises}\label{diophantine-equations:section-exercises}

\begin{exercise}\label{diophantine-equations:exercise-frobenius-three}
Show that $43$ is not of the form $6x + 9y + 20z$ with $x, y, z \in \NN$, but that every integer $n \geq 44$ is.
\end{exercise}

\begin{solution}
$43$ is not a multiple of $3$, so $z \neq 0$; $43 - 20 = 23$ and $43 - 40 = 3$ are not of the form $6x + 9y$ (the first is not a multiple of $3$, the second is too small). For $n \geq 44$: $44 = 6 \cdot 4 + 20$,
$45 = 9 \cdot 5$, $46 = 6 + 2 \cdot 20$, $47 = 3 \cdot 9 + 20$, $48 = 6 \cdot 8$, $49 = 9 + 2 \cdot 20$; these are six consecutive integers, and adding multiples of $6$ reaches every larger integer.
\end{solution}

\begin{exercise}\label{diophantine-equations:exercise-sixty}
Show that $60$ divides $xyz$ for every Pythagorean triple $(x, y, z)$.
\end{exercise}

\begin{solution}
It suffices to treat primitive triples $(m^2 - n^2, 2mn, m^2 + n^2)$. \emph{Divisibility by $4$}: one of $m$, $n$ is even, so $4 \mid 2mn$. \emph{By $3$}: if $3 \nmid mn$, then $m^2 \equiv n^2 \equiv 1 \pmod 3$
and $3 \mid m^2 - n^2$. \emph{By $5$}: if $5 \nmid mn$, then $m^2, n^2 \equiv \pm1 \pmod 5$, so $5$ divides $m^2 - n^2$ or $m^2 + n^2$.
\end{solution}

\begin{exercise}\label{diophantine-equations:exercise-markov-coprime}
Show that the entries of a Markov triple are pairwise coprime; in particular at most one of them is even.
\end{exercise}

\begin{solution}
If a prime $p$ divides two entries of a solution of \eqref{diophantine-equations:equation-markov}, it divides the square of the third, hence all three. This property is preserved by moves: if $p$ divides $x$, $y$,
$z$, it divides $3xy - z$. Conversely a move is an involution, so if $p$ divided all entries of some triple, it would divide all entries of $(1, 1, 1)$, by Theorem~\ref{diophantine-equations:theorem-markov}.
For parity: two even entries would have the common factor $2$.
\end{solution}

\begin{exercise}\label{diophantine-equations:exercise-powers-two-three}
Find all solutions in positive integers of $3^a - 2^b = 1$ and of $2^b - 3^a = 1$.
\end{exercise}

\begin{solution}
\emph{$3^a - 2^b = 1$.} If $b = 1$, $a = 1$. If $b \geq 2$, then $3^a \equiv 1 \pmod 4$, so $a = 2c$ is even and $2^b = (3^c - 1)(3^c + 1)$; both factors are powers of $2$ differing by $2$, so they are $2$ and
$4$, $c = 1$, $(a, b) = (2, 3)$. \emph{$2^b - 3^a = 1$.} Modulo $3$, $2^b \equiv 1$, so $b = 2c$ is even and $3^a = (2^c - 1)(2^c + 1)$; the factors are coprime powers of $3$ (their difference is $2$), so
$2^c - 1 = 1$, $c = 1$, $(a, b) = (1, 2)$. So the solutions are $3 - 2 = 1$, $9 - 8 = 1$ and $4 - 3 = 1$.
\end{solution}

\begin{exercise}\label{diophantine-equations:exercise-two-squares-conic}
Show that $x^2 + y^2 = 2z^2$ has infinitely many primitive solutions, and that $x^2 + y^2 = 3z^2$ has none other than $(0, 0, 0)$.
\end{exercise}

\begin{solution}
The rational points of $X^2 + Y^2 = 2$ found in Example~\ref{diophantine-equations:example-conics}(2) are infinitely many, and clearing denominators gives infinitely many primitive solutions, such as
$(1, 1, 1)$, $(1, 7, 5)$, $(7, 17, 13)$. For $3$: a nontrivial solution would give a rational point of $X^2 + Y^2 = 3$, which does not exist by Example~\ref{diophantine-equations:example-conics}(1).
\end{solution}

\begin{exercise}\label{diophantine-equations:exercise-not-square-diophantine}
Show that the set of positive integers that are not perfect squares is Diophantine.
\end{exercise}

\begin{solution}
$a$ is not a square if and only if $x^2 < a < (x + 1)^2$ for some $x \in \NN$, that is, if and only if there are $x, y, z \in \NN$ with $a = x^2 + 1 + y$ and $(x + 1)^2 = a + 1 + z$. By
Proposition~\ref{diophantine-equations:proposition-diophantine-closure} the conjunction of the two equations is the single equation $(a - x^2 - 1 - y)^2 + (x^2 + 2x - a - z)^2 = 0$.
\end{solution}
